\section{Clases y métodos implementados por cada subequipo}

\textbf{Cortez Vega y Lujan De la Rosa:}

\begin{enumerate}
  \item \textbf{Módulo Algoritmo PRIMERO} (\texttt{algoritmo\_primeros.py}): Implementa la
    función \texttt{calcular\_primeros}, que calcula los conjuntos PRIMERO para todos los
    símbolos de la gramática mediante un algoritmo iterativo de punto fijo. La función recibe
    el símbolo de consulta, el diccionario de producciones, un diccionario mutable de conjuntos
    \texttt{primeros} y la lista de terminales. Inicializa cada terminal con su propio conjunto
    y cada no terminal con el conjunto vacío; luego itera hasta que no haya cambios.

  \item \textbf{Módulo Algoritmo SIGUIENTE} (\texttt{algoritmo\_siguientes.py}): Implementa
    dos funciones:
    \begin{enumerate}[label=\alph*.]
      \item \textbf{\texttt{\_primero\_secuencia}}: función auxiliar que calcula
        $\text{PRIMERO}$ de una secuencia de símbolos e indica si toda la secuencia puede
        derivar en $\varepsilon$.
      \item \textbf{\texttt{calcular\_siguientes}}: función principal que aplica el algoritmo
        de SIGUIENTE mediante iteración de punto fijo. Inicializa el símbolo inicial con
        $\{\$\}$ y, para cada producción $A \to \alpha\, B\, \beta$, propaga
        $\text{PRIMERO}(\beta)$ y, cuando corresponde, $\text{SIGUIENTE}(A)$ hacia
        $\text{SIGUIENTE}(B)$.
    \end{enumerate}
\end{enumerate}

\vspace{0.8cm}

\textbf{Arias Mendoza y Vasquez Moreno:}

\begin{enumerate}
  \item \textbf{Módulo Gramática} (\texttt{gramatica.py}): Define la clase \texttt{Gramatica},
    que representa una gramática libre de contexto. Sus atributos son:
    \texttt{terminales}, \texttt{no\_terminales}, \texttt{simbolo\_inicial} y
    \texttt{producciones} (diccionario de cabeza a lista de cuerpos). El método
    \texttt{agregar\_produccion} inserta una nueva regla de producción al diccionario.

  \item \textbf{Módulo Lector de Archivo} (\texttt{lector\_archivo.py}): Implementa la función
    \texttt{cargar\_gramatica\_desde\_archivo}, que lee un archivo \texttt{.txt} codificado en
    UTF-8 y construye un objeto \texttt{Gramatica}. Soporta dos formatos:
    \begin{enumerate}[label=\alph*.]
      \item \textbf{Formato simple}: solo producciones; terminales y no terminales se infieren
        automáticamente — cualquier símbolo que encabece una producción es no terminal, el
        resto son terminales.
      \item \textbf{Formato con encabezados}: las dos primeras líneas declaran explícitamente
        los terminales y no terminales; las producciones inician en la tercera línea.
    \end{enumerate}
    Acepta producciones vacías escritas como $\varepsilon$, \texttt{e} o \texttt{epsilon}
    (sin distinción de mayúsculas) y lanza \texttt{ValueError} ante cualquier producción
    malformada.
\end{enumerate}

\vspace{0.8cm}

\newpage

\textbf{Carrasco Pacheco y Nicolas Vasquez:}

\begin{enumerate}
  \item \textbf{Módulo de Estilos} (\texttt{estilos.py}): Define la paleta de colores
    (\texttt{COLORES}), la función \texttt{configurar\_estilos} que aplica los temas
    \texttt{ttk} globalmente, y la función \texttt{crear\_superficie} que genera marcos con
    un diseño visual uniforme para todas las secciones de la interfaz.

  \item \textbf{Módulo Tabla Base} (\texttt{tabla\_base.py}): Widget \texttt{Treeview}
    reutilizable que incluye barras de desplazamiento vertical y horizontal. Permite
    configurar las columnas, sus encabezados, anchos y alineación, y es utilizado como
    base por los demás módulos de presentación.

  \item \textbf{Módulo Tabla Conjuntos} (\texttt{tabla\_conjuntos.py}): Especializa la tabla
    base para mostrar los conjuntos PRIMERO o SIGUIENTE, presentando un no terminal por
    fila y su conjunto correspondiente como cadena de símbolos separados por comas.

  \item \textbf{Módulo Ventana de Resultados} (\texttt{ventana\_primeros\_siguientes.py}):
    Implementa la clase \texttt{VentanaPrimerosSiguientes}, una ventana secundaria con dos
    pestañas: una para los conjuntos PRIMERO y otra para los conjuntos SIGUIENTE. Cada pestaña
    muestra una \texttt{TablaConjuntos} con los resultados del análisis.

  \item \textbf{Módulo Interfaz Principal} (\texttt{interfaz.py}): Implementa la clase
    \texttt{InterfazCompilador}, punto de entrada de la aplicación. Proporciona:
    \begin{itemize}
      \item Carga de archivos \texttt{.txt} mediante diálogo del sistema operativo.
      \item Editor de solo lectura con numeración de líneas para inspeccionar la gramática.
      \item Botón \textbf{Analizar} que invoca los módulos de PRIMERO y SIGUIENTE y abre la
            ventana de resultados.
      \item Botón \textbf{Limpiar} para reiniciar el estado de la aplicación.
      \item Barra de estado con mensajes informativos sobre el progreso del análisis.
    \end{itemize}
\end{enumerate}
